\section{Emerging and Specialized Protocols}
\label{sec:emerging-protocols}

This section applies the rubric to the protocols whose
2026-Q1 maturity is at the academic-paper, arXiv-draft, or
EIP-draft stage rather than at the production-deployment
stage of the protocols in \S\ref{sec:industry-protocols}.
The lower deployment maturity does not imply lower
significance: ACNBP's capability-binding ceremony and LOKA's
verifiable-credential framework address gaps that the
industry-standard protocols have left open, and the five
commerce-settlement protocols introduce an architectural
layer that the prior surveys did not cover.

\subsection{LOKA}
\label{sec:emerging:loka}

The LOKA Protocol~\cite{spec_loka} is the most
ambitious peer-messaging proposal of the surveyed set in
the sense that it commits, at the specification level, to
the strongest version of every axis where the
industry-standard protocols make pragmatic concessions.
LOKA proposes a Semantic Discovery Fabric (Axis 1
score = 3), DIDs with Ed25519 keys and post-quantum
Dilithium upgrades (Axis 2 score = 3), VCs as scoped
capability attestations combined with smart-contract
lifecycle policies (Axis 3 score = 2), an
intent-centric peer-to-peer async model with continuous
checkpointing and recovery (Axis 6 score = 3), a Polyglot
Intent Engine targeting FIPA-ACL, A2A, and Open Voice
interoperability (Axis 9 score = 3 aspirationally), and
immutable blockchain-backed audit trails with real-time
ethical auditing (Axis 10 observability = 3).

The authors themselves note in a footnote that LOKA is a
conceptual research framework; the cited specification
declares mechanisms but does not yet ship reference
implementations of most of them. The Axis 4 transport
tuple records this: \texttt{transports = unspecified},
\texttt{streaming = unspecified}, \texttt{async =
intent-centric peer-to-peer (unspecified)}~\cite{spec_loka}.
The strength of LOKA in the rubric matrix is therefore
that it occupies the upper-right corner on most axes; the
weakness is that the implementation-maturity context that
\S\ref{sec:rubric:limits} flags applies in full. A
reader weighting the survey by source tier should treat
LOKA's scores as aspirational claims about a possible
future protocol rather than as verifiable claims about a
deployed one.

The Axis 8 score of 6 of 12 reflects an explicit
threat-modeling pass in the LOKA paper that addresses
DID-and-Dilithium authentication, timestamp-and-signed-intent
replay, blockchain audit, post-quantum MITM,
decentralized-SSI registry, and Contextual Ethics privilege
confusion. Four further classes are treated but not settled,
so they score partial and do not enter the count: the DECP
prompt-injection filter, VC-and-contract capability scoping,
delegation token binding, and the privacy side of Contextual
Ethics. Supply-chain attacks and DoS are unaddressed.

\subsection{Agent Capability Negotiation and Binding Protocol (ACNBP)}
\label{sec:emerging:acnbp}

ACNBP~\cite{spec_acnbp} is the most security-rigorous of the
emerging protocols in the rubric matrix. Its Axis 8 score of
10 of 12 is the highest of any protocol surveyed, the
direct consequence of an explicit MAESTRO 7-layer threat
analysis in the specification that addresses zero-knowledge
proof capability attestations, signature plus HMAC
integrity, rate-limit-and-proof-of-work DoS mitigation,
nonce-and-timestamp-and-sequence replay defense, a
\texttt{protocolExtension} whitelist for supply-chain
attacks, X.509 with OCSP against agent impersonation, the
Capability Binding ceremony against delegation-token theft,
signed Agent Network Resource Identifiers with BFT for
registry integrity, foundation-model-layer prompt injection
mitigation, and cryptographically protected audit logs. The
two classes that score partial, and so do not enter the
count, are side-channel privacy through ring signatures and
cross-agent privilege confusion through agent quarantine.

The protocol's structural contribution is the ten-step
capability negotiation and binding flow. Two agents that
have located each other through the ANS directory
exchange capability descriptors, run a semantic
consistency check, propose binding terms, negotiate
modifications, and conclude with a mutually-signed
Capability Binding (the central primitive defined in
Definition 10 of the specification). The bound capability
is then executable as an artifact-producing operation
whose execution produces a cryptographically attested
receipt. The Axis 5 score of 3 (10-step state machine
with Protocol State as a first-class object) follows
directly from this structure.

Authorization is where that structure stops, and the
boundary is worth stating precisely because it is easy to
misread. Definition 10 binds two parties: a requester and
a provider, with capabilities, terms, and mutual
signatures. That is one hop, and it is bound more
rigorously than any other peer-messaging specification
binds one hop. It is not a chain. The ACNBP text never
uses the words delegation, authorization, on behalf,
principal, or sub-agent, and the nearest multi-agent
construct, the Multi-Agent Orchestration paragraph of
\S IV.C, describes capability composition rather than
transitive authority. The Axis 3 level-3 anchor
(\S\ref{sec:rubric:axis3}) asks whether a downstream
agent can independently verify an intermediate agent's
grant. Nothing in ACNBP answers that, so the Axis 3
score is 2, which is the anchor for a structured
capability with explicit scope and no standardized
chain-of-custody format. A lead author
confirmed the reading in correspondence: ACNBP is a
two-party protocol for service negotiation and binding,
and multi-agent delegation is outside its scope. The
consequence for the field is in
\S\ref{sec:comparison:derivation}, and it is the
strongest form of the survey's central finding: no
scored protocol reaches level 3, including the most
security-rigorous one.

Identity is the one axis where ACNBP's own text is
narrower than its authors' work. The specification
authenticates agents with X.509 certificates, OCSP
revocation, and signed envelopes, which is Axis 2 level 2,
and it never mentions decentralized identifiers or
verifiable credentials. The same authors do treat
decentralized identity, in a companion framework built on
DIDs and verifiable credentials with a naming service and
zero-knowledge attribute
disclosure~\cite{huang_zerotrust_identity}. ACNBP cites
that framework as related work and does not require it, so
the Axis 2 score records what ACNBP itself specifies. The
survey scores specifications, not the composition a reader
could assemble from a research group's output
(\S\ref{sec:rubric:method}). Were the companion framework
folded in normatively, ACNBP would be a candidate for
Axis 2 level 3.

The principal limitation at the cutoff is governance: the
specification is authored by individuals affiliated with
DistributedApps.ai, Cisco, OWASP/AWS, and Intuit, but is
not yet under a formal SDO. The Axis 10 governance
sub-score therefore reflects ``individual authors''
rather than a foundation-or-W3C tier. The technical
maturity is higher than the institutional maturity, which
is the inverse of the situation for A2A and MCP.

\subsection{The agentic commerce stack}
\label{sec:emerging:commerce}

The fourth and structurally distinct group of protocols is
the agentic commerce stack: AP2~\cite{spec_ap2}, Visa
Trusted Agent Protocol~\cite{spec_visa_tap}, Mastercard
Agent Pay, ERC-8004~\cite{spec_erc8004}, and
UCP~\cite{spec_ucp}. The first four are a co-stack
targeting different settlement substrates (card networks,
on-chain) rather than competitors. UCP is a different kind
of thing and is treated separately below. The taxonomy of \S\ref{sec:taxonomy} places
them on the third architectural layer
(commerce-settlement) for a reason: they do not specify
how agents communicate with each other in the
peer-messaging sense, and they do not specify how an
agent invokes external capabilities in the tool-binding
sense. They specify what is necessary for one agent to
prove, at the moment of payment, that it is acting within
the authority granted by a human principal. Neither the
peer-messaging nor the tool-binding layers
address that concern.

The architectural primitive shared by all four is the
\emph{intent-mandate-receipt envelope}:

\begin{itemize}
  \item The \textbf{intent} expresses what the agent
    proposes to do (e.g., purchase item X for at most
    price Y from merchant Z within timeframe T). It is
    structured and machine-verifiable rather than
    free-text.
  \item The \textbf{mandate} is the principal's signed
    authorization for the agent to act within the
    intent's scope. The mandate is bounded in scope
    (specific intent classes, specific merchants),
    bounded in time (expires), and revocable.
  \item The \textbf{receipt} is the cryptographically
    attested record of execution, signed by the
    settlement authority (card network, payment
    processor, or on-chain settlement contract), that
    binds the execution to the mandate and the intent.
\end{itemize}

Four of the five protocols share this envelope but
differentiate on steward, settlement substrate, mandate
encoding, and on-chain versus off-chain anchor;
Table~\ref{tab:commerce-stack} summarizes the contrast.

AP2 is the PayPal-and-Google-led specification of the
envelope for general-purpose agentic
commerce~\cite{spec_ap2}, with reference implementations
targeting PayPal's payment rails. Visa TAP~\cite{spec_visa_tap}
binds the envelope to the Visa card-network settlement
substrate, with explicit cryptographic primitives that
align with EMVCo-compatible authorization flows.
Mastercard Agent Pay occupies the same architectural slot
for the Mastercard network. ERC-8004~\cite{spec_erc8004}
binds the envelope to on-chain settlement on Ethereum and
EVM-compatible chains, with the mandate expressed as a
smart contract and the receipt as a transaction record.
The on-chain route inherits a standards and trust-boundary
literature that predates the agentic-commerce
stack~\cite{agents_blockchain}. That work treats the
execution model and the limits of on-chain trust in detail.

\paragraph{UCP, the orchestration layer above the
envelope.} The four protocols above answer one question:
how a payment is authorized and settled. UCP answers a
different one, and that is why it sits in this section
without sharing the envelope. It specifies the commerce
transaction that surrounds the payment. The surveyed
version, \texttt{2026-01-11}, defines three standard
capabilities: Checkout, which covers cart building and tax
calculation, Order, which pushes lifecycle updates, and
Identity Linking~\cite{spec_ucp}. Settlement is not among
them. UCP binds no substrate of its own. Payment
authorization is delegated to an optional AP2 Mandates
extension (\texttt{dev.ucp.shopping.ap2\_mandate}). Once that
extension enters the negotiated capability set, the session is
\emph{Security Locked}. Neither party may then revert to an
unprotected checkout.

That relationship is the reason UCP matters to this
survey beyond its own scope. It is a working instance of
the layered composition that \S\ref{sec:comparison:stack}
argues for, built by parties who had no stake in the
argument. UCP consumes AP2 at the settlement layer. It
also declares four transport bindings. They are REST via
OpenAPI, MCP via OpenRPC, A2A with the Agent Card URL as the
endpoint, and an embedded binding for merchant-hosted user
interfaces. Two of the four are protocols this survey scores.
A commerce-layer protocol that treats A2A and MCP as
interchangeable carriers is evidence that the three
architectural layers of \S\ref{sec:taxonomy:layers}
compose in practice, and not only on paper. We contacted
Shopify about this reading, and its specification authors
confirmed it.

UCP is also the one commerce-layer protocol with
machinery on the axes where the other four are silent. It
specifies discovery: businesses publish a profile at
\texttt{/.well-known/ucp}, and platforms declare their own
profile URI on every request through a \texttt{UCP-Agent}
header. Capability names are reverse-domain strings, so
naming authority follows domain ownership. That puts
\texttt{dev.ucp.*} under the governing body and leaves
vendor extensions under \texttt{com.\{vendor\}.*}. It
specifies authorization: an OAuth 2.0 authorization-code
flow for identity linking, with per-capability scopes such
as \texttt{ucp:scopes:checkout\_session} and a
\texttt{/.well-known/oauth-authorization-server} document.
Authentication is specified in two places rather than one.
Order webhooks are signed with a detached JWT (RFC 7797)
over the request body. The \texttt{kid} header selects a key
from a \texttt{signing\_keys} JWK array in the sender's
profile.
The AP2 extension signs checkout terms with JWS detached
content and requires ES256, ES384, or ES512.

Request signing and identity linking answer two different
questions, and reading them as alternatives is a mistake this
survey initially made. The same exchange with Shopify
corrected it. A signature establishes which
platform sent the payload, so it is agent identity. The OAuth
flow establishes which buyer the platform is acting for, so it
is buyer identity. Neither substitutes for the other, and a
deployment chooses how much of each it wants. Three shapes
follow. A business may welcome any agent and accept guest
buyers. It may instead accept only a closed list of signed
agents. It may require buyer sign-in as well, which is the
common combination in business-to-business flows. The variance
is deliberate, so the specification sets a floor and leaves the
ceiling to the deployment.

Three limits are worth stating plainly, and the first is
narrower than it first appears. The authentication requirement
is asymmetric. Business-to-platform webhooks \textbf{MUST} be
signed, while platform-to-business requests only
\textbf{SHOULD} be authenticated, with a bearer token as the
example given. The asymmetry has a reason. A webhook receiver
exposes an open endpoint, so the pushed payload has to carry
its own proof of origin, which is what a detached JWT supplies.
A platform calling a business calls a known endpoint over TLS,
where that argument does not hold. Shopify's authors
confirmed this reading and added that the \textbf{SHOULD} is a floor
rather than a ceiling, since some businesses do enforce it as a
\textbf{MUST}. So the limit is not that the weaker direction
goes unprotected in practice. It is that the protocol makes no
guarantee either way, so a receiver cannot infer from UCP
conformance alone whether the platform calling it authenticated
at all. Second, request signing is scoped to the
webhook and the AP2 extension rather than defined once for the
protocol. A deployment that negotiates neither has
transport security and a bearer token. That is Axis 2 level 1
behaviour underneath a specification that reads as
stronger. Third, delegation stops where the other protocols
stop. Identity linking scopes what one platform may do on a
user's behalf at one business. It carries no chain, so a
platform acting through a second platform cannot prove the
provenance of its authority. The strong delegation primitive
is available only through the AP2 extension, and only for the
payment step. UCP is therefore a fifth independent instance
of Finding~3 rather than a counterexample to it. It is also
the instance drawn from the protocol with the most commercial
weight behind it.

Governance is the other place UCP does not fit the existing
patterns. There is no formal standards body. A Tech Council
holds technical authority over the specification and a
Governing Council holds authority over the project. Both are
weighted to the two founding organisations. The Tech Council
has twelve votes. Eight are reserved, four each to Google and
Shopify. The other four are elected every six months from any
organisation. The Governing Council has three seats: one
Google, one Shopify, and one elected annually. The elected
seat is unfilled at the surveyed version, and Google holds its
proxy. Google is also custodian of the \texttt{ucp.dev}
domain. Domain Working Groups may be chartered by three or
more organisations, and every artifact they produce needs Tech
Council approval. The specification is Apache 2.0 and
contributions require a signed contributor licence.

This is a fourth governance pattern, distinct from the three
that \S\ref{sec:governance} identifies. It is neither a
foundation, nor a single vendor, nor an academic group with no
steward. It is a two-vendor partnership with a stated
open-ecosystem mandate and a structural majority reserved to
the founders. The gap between the two is the thing to watch.
The governance document asks members to act ``agnostic to
interests of the companies they represent''. It also gives
those companies two thirds of the votes that decide the
specification. The elected seats and the published meeting
minutes are the mechanisms that would make the mandate real.
Neither had produced an observable record at the surveyed
version, so how the model behaves under disagreement between
Google and Shopify is not yet knowable.

The rubric applies awkwardly to the commerce-settlement
layer because several of the peer-messaging axes are
structurally inapplicable. Axis 5 (task and conversation
model) and Axis 7 (multi-modal content handling) have no
counterpart in a commerce-settlement protocol whose unit
of interaction is the intent-mandate-receipt triple
rather than a multi-message conversation. Axis 3
(authorization and delegation) is the axis on which the
commerce stack defines its own ceiling: the
mandate-as-signed-capability-with-bound-intent
architecture is the most rigorous delegation model in any
of the surveyed protocols, and the cross-protocol
interoperability work underway in the AAIF aims to make
the same mandate format consumable across A2A's
peer-messaging layer and ANP's identity layer.

The principal open issue at the cutoff is interoperability
across the four settlement protocols: a single agent
acting on behalf of a principal would currently need to
bind to AP2, Visa TAP, Mastercard Agent Pay, and
ERC-8004 separately if it wanted to transact across all
four settlement substrates. The AAIF's Agentic Commerce
Working Group~\cite{lf_aaif} has begun harmonizing the
envelope schema; the cutoff state is that the four
protocols share the architectural pattern but not the
wire format.

UCP shows one way that gap gets closed, and it is not the
way the harmonization effort assumes. Rather than
converging the four wire formats, UCP fixes the
transaction layer above them and treats the settlement
mandate as a negotiated extension. An agent then binds to
one orchestration protocol and inherits whichever
settlement substrate the extension names. At the cutoff
the only such extension is AP2, so the approach is
demonstrated on one substrate and not yet on four. Which
route prevails, a harmonized envelope or an orchestration
layer with pluggable mandates, is an open question that
\S\ref{sec:openproblems} returns to.

\subsection{Semantic-driven proposals}
\label{sec:emerging:semantic}

A small body of work~\cite{semantic_driven,
beyond_message_passing} proposes augmenting or replacing
the LLM-era protocols with semantic-first designs that
re-introduce, in modernized form, the ontology commitments
of FIPA-ACL. The \emph{Beyond Message Passing} line of
work~\cite{beyond_message_passing} argues that the
free-text Skill descriptions and tool descriptions of the
industry-standard protocols (Axis 9 = 1) are
fundamentally insufficient for cross-organizational
interoperability and proposes a vocabulary-and-matcher
architecture. The \emph{Semantic-Driven AI Agent
Communications} line~\cite{semantic_driven} takes a
complementary tack, proposing intent-and-capability
ontologies that adapt FIPA-SL to a JSON-LD substrate.

These proposals are at a pre-specification stage as of
the cutoff and do not yet have a place in the rubric
matrix. Section~\ref{sec:semantic} treats the underlying
gap that they address.
